\paragraph{SSH Client Authentication}

Password-based authentication of SSH clients has been studied previously by Owens and Matthews, and Song et al.~\cite{USENIX:SonWagTia01,owens2008study}. Javed and Paxson~\cite{CCS:JavPax13} statistically measured and evaluated distributed online dictionary attacks against password-based authentication in SSH. To identify client authentication efforts using a stolen SSH client password, Piet et al.~\cite{USENIX:PSPW23} used keystroke authentication. Park et al.~\cite{park2021network} analyzed router log data to detect SSH brute force attacks. SSH honeypots were used by Melese and Avadhani, Koniaris et al., and Wu et al.~\cite{6624967,melese2016honeypot,Wu2020MiningTI} to study brute-force online dictionary attacks. In this paper, we study the security of SSH client authentication with digital signatures.

\paragraph{SSH Client Keys}

In 2015, Cartwright-Cox~\cite{cox2015} scraped 1.2M keys of GitHub users, found 49~RSA keys shorter than 1000~bits, and reported some users whose keys were generated on a Debian system with weak randomness. In 2023, Böck~\cite{EPRINT:Bock23} tested the Fermat attack against this data set and found no vulnerable keys. In 2016, Barbulescu et al.~\cite{barbulescu2016weakrsakeys} analyzed weak RSA keys on the internet, including 3M RSA keys from GitHub users. They found one 512-bit key, but could not find any common factors using batch GCD among the other keys. Also, in 2015, Cryptosense~\cite{cryptosense2015} found small factors in GitHub user keys and could factor~12 keys with 2048~bits using GMP\nobreakdashes-ECM by Zimmermann.\footnote{\url{https://gitlab.inria.fr/zimmerma/ecm}, accessed: 2025-04-14} In 2018, Amiet and Romailler~\cite{amiet2018} analyzed 4.7M keys from GitHub and 1.7M keys from GitLab for ROCA vulnerabilities and common factors using batch GCD. They found vulnerable keys in their dataset, but did not separate results by data source. In this work, we collect keys from GitHub, GitLab, and Launchpad. After reviewing the related work and relevant standards, we identified a broad range of security issues. Based on these issues, we implemented---to the best of our knowledge---the most comprehensive suite of security tests for the detection of potential security issues and vulnerabilities.

\paragraph{SSH Signature-based Authentication}

SSH uses the \pkcsone signature scheme for authentication with RSA keys. Yahyazadeh et al.~\cite{CCS:YCLHDI21} formalized the \pkcsone signature specification and developed a black-box differential testing framework targeting this signature scheme. After analyzing 45 libraries, they identified various conformance issues and six libraries susceptible to Bleichenbacher's signature forgery attacks~\cite{bleichenbacher2006}. These attacks allow an attacker to forge RSA signatures if the verifier does not strictly check the \pkcsone format of the verified signature, and the RSA public key exponent is small (e.g., $e=3$), which we did not observe in the collected data sets. In our work, we concentrated on the security analysis related to signature generation. Signature verification evaluations, such as those presented by Yahyazadeh et al.~\cite{CCS:YCLHDI21}, could extend in our framework.

Ryan et al. recovered private RSA keys of SSH servers from faulty signatures in~\cite{CCS:RHSH23}. Breitner and Heninger studied ECDSA signatures in cryptocurrencies and SSH, and were able to extract private keys due to biased nonces~\cite{FC:BreHen19}. By collecting public SSH server keys from the internet, Heninger et al. could compute the private keys of 0.03\% of SSH servers using RSA and 1.03\% of SSH servers using DSA~\cite{USENIX:HDWH12}. In this paper, we use techniques similar to Heninger et al. to analyze the security of SSH \emph{client} keys.

\paragraph{Deterministic Signature Schemes}

Cao et al. and Poddebniak et al.~\cite{EUROSP:PSSLR18,RSA:CSCCFW22} showed that deterministic variants of ECDSA and EdDSA are vulnerable to \emph{fault attacks}, where invalid signatures are needed. Please note that our attack on \putty is not a fault attack but uses \emph{valid} signatures. We used the algorithm described by Albrecht and Heninger~\cite{EC:AlbHen21} to compute the private key.

\paragraph{Attacks on the SSH Protocol}

In 2009, Albrecht, Paterson, and Watson~\cite{SP:AlbPatWat09} attacked the binary packet protocol of SSH by exploiting the encrypted length field. By measuring the number of ciphertext bytes accepted by the server from the network, they derived a decryption oracle for some bits of an encrypted block. Paterson and Watson~\cite{EC:PatWat10} showed a variation of this attack. Wei Dai~\cite{weidai2002} described a theoretical attack against SSH employing CBC-mode ciphers, while Bhargavan and Leurent~\cite{NDSS:BhaLeu16} presented a handshake transcript collision attack using \shaone. Bäumer et al.~\cite{USENIX:BauBriSch24} showed a prefix truncation attack on the ciphertext for some cipher modes using sequence number manipulation, despite formal security proofs for SSH~\cite{CCS:BelKohNam02,EC:PatWat10,CCS:ADHP16}. Fiterau-Brostean et al.~\cite{fiteruau2017model} learned Mealy state machines of three SSH server implementations; no weaknesses were discovered.
